\section{Mechanism and Robustness}
\label{sec:robustness}

This section separates the economic mechanism from the functional forms used in the baseline. The results below are restricted robustness statements; they do not establish a theorem for arbitrary continuous portfolio games, matching technologies, or numbers of regions and activities.

\subsection{A binary switching representation}

Consider a binary comparison between a duplicated configuration and a configuration in which one jurisdiction reallocates its policy toward a complementary activity. Let $\Delta>0$ be the direct local surplus sacrificed by the reallocating jurisdiction, let $G>0$ be the total complementary surplus created by the switch, and let $\lambda\in(0,1)$ be the fraction of that surplus captured by the reallocating jurisdiction.

\begin{proposition}[Incomplete-capture switching wedge]
\label{prop:switching-lemma}
The coordinated switch is desirable if and only if
\begin{equation}
G>\Delta,
\label{eq:planner-G}
\end{equation}
whereas the reallocating local government switches if and only if
\begin{equation}
\lambda G>\Delta.
\label{eq:local-G}
\end{equation}
Thus the coordinated and decentralized switching thresholds are
\begin{equation}
G^P=\Delta,
\qquad
G^N=\frac{\Delta}{\lambda},
\label{eq:G-thresholds}
\end{equation}
and $G^P<G^N$ whenever $0<\lambda<1$.
\end{proposition}

The baseline is a special case with $G=A$ and $\lambda=1-\alpha$ for a region switching from $U$ to $D$ while the other remains at $U$. Proposition \ref{prop:switching-lemma} is a sufficient-condition theorem for a finite configuration switch; it is not a claim about arbitrary continuous portfolio games.

If $G=G(\tau)$ is continuous and strictly increasing in an integration index $\tau$, $\lambda$ is constant, and both $\Delta$ and $\Delta/\lambda$ lie in the range of $G$, unique thresholds satisfy
\begin{equation}
G(\tau^P)=\Delta,
\qquad
G(\tau^N)=\frac{\Delta}{\lambda},
\end{equation}
so that $\tau^P<\tau^N$. The crossing assumptions matter. If the second threshold does not exist, decentralized reallocation need not occur at any finite level of integration.

The same caution applies when the capture share is endogenous. If $\lambda=\lambda(\tau)$, define $H(\tau)=\lambda(\tau)G(\tau)$. The planner compares $G(\tau)$ with $\Delta$, whereas the local government compares $H(\tau)$ with $\Delta$. Incomplete capture alone does not imply a finite local threshold. For example, with $G(\tau)=1+\tau$, $\lambda(\tau)=(1+\tau)^{-2}$, and $\Delta=1$, the planner prefers reallocation for every $\tau>0$, while $H(\tau)=1/(1+\tau)\leq1$, so the local government never switches.

\subsection{Production beyond Cobb--Douglas}

The incidence logic does not require Cobb--Douglas production. Consider
\begin{equation}
Y=A F(u,d),
\label{eq:general-production}
\end{equation}
where $F$ is differentiable, concave, homogeneous of degree one, and satisfies $F_u(1,1)>0$ and $F_d(1,1)>0$. Under competitive factor pricing, Euler's theorem gives
\begin{equation}
F(1,1)=F_u(1,1)+F_d(1,1).
\end{equation}
For the same binary assignment comparison, the coordinated threshold is
\begin{equation}
A^P_F=\frac{\Delta}{F(1,1)},
\end{equation}
while the $D$ jurisdiction's local threshold is
\begin{equation}
A^N_F=\frac{\Delta}{F_d(1,1)}.
\end{equation}
Since $F_u(1,1)>0$, $F(1,1)>F_d(1,1)$ and therefore $A^P_F<A^N_F$. The wedge is again generated by incomplete local capture of total complementary value.

As a check, for the concave CES technology
\begin{equation}
F(u,d)=\left[\omega u^\rho+(1-\omega)d^\rho\right]^{1/\rho},
\qquad 0<\omega<1,
\end{equation}
within its admissible concave parameter domain, $F(1,1)=1$ and $F_d(1,1)=1-\omega$. Hence the local threshold is $\Delta/(1-\omega)$. This exercise removes dependence on the Cobb--Douglas exponent while retaining the same incidence mechanism.

\subsection{Alternative matching technology}

The baseline uses random matching. A qualitatively different capacity-matching technology pairs as many available complementary projects as possible. Let
\begin{equation}
m_{iU,jD}=\min\{x_i,1-x_j\},
\qquad
m_{iD,jU}=\min\{1-x_i,x_j\},
\end{equation}
so total complementary matches are
\begin{equation}
M(x_1,x_2)
=m_{1U,2D}+m_{1D,2U}
=1-|x_1+x_2-1|.
\label{eq:capacity-matching}
\end{equation}
Let the region supplying $D$ capture fraction $\lambda\in(0,1)$ of a match's surplus and the region supplying $U$ capture $1-\lambda$. Region $i$'s complementary income is therefore
\begin{equation}
(1-\lambda)A m_{iU,jD}+\lambda A m_{iD,jU}.
\label{eq:capacity-local-income}
\end{equation}
For any rival portfolio, region $i$'s payoff is increasing in $x_i$ on the first piece of the matching function. On the second piece its slope is $\Delta-\lambda A$. Hence, if $A<\Delta/\lambda$, $x_i=1$ is a strict best response to every $x_j$, so $(1,1)$ is the unique decentralized equilibrium.

The coordinated objective, apart from constants, is
\begin{equation}
\Delta(x_1+x_2)+A M(x_1,x_2).
\end{equation}
If $A>\Delta$, its optimum set is
\begin{equation}
x_1+x_2=1.
\label{eq:capacity-optima}
\end{equation}
Thus for
\begin{equation}
\Delta<A<\frac{\Delta}{\lambda},
\end{equation}
decentralization uniquely duplicates the high-return activity while coordination reallocates policy away from duplication. The threshold mechanism therefore survives this second matching technology, but the exact geographic form of the coordinated allocation changes. In particular, $(1/2,1/2)$ can be coordinated-optimal under \eqref{eq:capacity-matching}. The robust conclusion is reallocation away from excessive duplication, not complete regional specialization.
